\chapter{Hash Functions}
\label{chap:hash-functions}
\raggedbottom

\section*{Author's Intuition}

The original draft sought a way to use a digest to ``express \(x\) without
revealing it.'' That is an important question, especially in adversarial
cooperation, but a bare hash does not answer it in general. This chapter keeps
the intuition visible and draws the boundary precisely: a digest is a public,
deterministic fingerprint of exact bytes. Whether it hides anything depends on
the input distribution and on a larger construction.

\begin{abstract}
This appendix introduces cryptographic hash functions as deterministic maps
from byte strings to fixed-length digests. It separates preimage,
second-preimage, and collision resistance; derives the familiar generic work
factors only in an idealized model; and distinguishes an observed avalanche
effect from a security proof. It then explains why an unkeyed digest is not
encryption, anonymity, a message-authentication code, a signature, a
commitment, password hashing, or a key-derivation function. The C companion
uses the repository's sole cryptographic backend, libsodium, to compute
BLAKE2b-256 test vectors. MD5 appears only as a historical collision example.
\end{abstract}

\section{A Function on Bytes}

For this appendix, an \emph{unkeyed \(n\)-bit hash function} is an efficiently
computable deterministic function
\[
    H : \{0,1\}^{*} \longrightarrow \{0,1\}^{n}.
\]
It accepts an arbitrary finite byte string and returns an \(n\)-bit digest.
The same bytes always produce the same digest. Different bytes are allowed to
produce the same digest: because the domain is larger than the finite range,
collisions necessarily exist. The cryptographic objective is not to abolish
collisions, but to make specified adversarial searches computationally
difficult.

This functional definition does not say that \(H\) is ``one way.'' One-wayness
is a security property defined by an experiment, an input or challenge
distribution, an adversary's resources, and its success probability. Formal
treatments also distinguish fixed functions from keyed hash-function families
and give several non-equivalent versions of the basic notions
\cite{rogaway2004hashbasics}. The experiments below are a deliberately compact
teaching model, not a replacement for those definitions.

\subsection{Bytes, not meanings}

The hash function sees neither words nor objects. It sees bytes. These inputs
are distinct:

\begin{verbatim}
61 62 63          ASCII "abc"
61 62 63 0a       ASCII "abc" followed by a newline
61 62 63 00       ASCII "abc" followed by a NUL byte
\end{verbatim}

Text encoding, normalization, field order, lengths, and terminators therefore
belong to the protocol specification. A digest cannot repair an ambiguous or
inconsistent encoding.

\section{Three Different Resistance Experiments}

Let \(H\) be fixed for the following informal experiments. A complete formal
statement would additionally specify how challenges are sampled and bound the
adversary's time, memory, queries, and success probability.

\subsection{Preimage resistance}

A challenge digest \(y\) is generated according to a stated experiment. The
adversary receives \(y\) and succeeds if it returns any \(x'\) such that
\[
    H(x') = y.
\]
The notation \(H^{-1}(y)\) can be misleading: a compressing hash is generally
many-to-one and has no unique inverse. The question is whether the adversary
can find one member of the preimage set under the stated challenge model.

\subsection{Second-preimage resistance}

The adversary receives a particular input \(x\) and succeeds if it returns a
different input \(x'\) such that
\[
    x' \ne x
    \qquad\text{and}\qquad
    H(x') = H(x).
\]
Here the first message is fixed by the experiment; the adversary does not get
to choose both messages freely.

\subsection{Collision resistance}

The adversary succeeds if it returns any pair \((x,x')\) satisfying
\[
    x \ne x'
    \qquad\text{and}\qquad
    H(x) = H(x').
\]
This is a different task because the adversary chooses both inputs. A practical
collision attack does not automatically give an efficient preimage attack for
a separately specified target digest.

The three names are related intuitions, not interchangeable labels. Even basic
implications depend on the chosen formalization and the amount of compression
\cite{rogaway2004hashbasics}.

\section{Idealized Generic Work Factors}

Suppose only for this calculation that distinct queries behave like independent
uniform samples from an \(n\)-bit range, and consider a classical generic
attacker.

For a fixed target \(y\), one trial succeeds with probability \(2^{-n}\). After
\(q\) independent trials, the success probability is approximately
\[
    \frac{q}{2^n}
\]
while that ratio remains small. A constant success probability therefore
appears at a scale near \(2^n\) trials.

Among \(q\) sampled outputs there are \(q(q-1)/2\) pairs. The expected number
of colliding pairs is approximately
\[
    \frac{q(q-1)}{2^{n+1}}.
\]
This becomes constant when \(q\) is on the scale of \(2^{n/2}\), the birthday
scale.

These are generic query estimates in an idealized model. They are not
universal Big-O guarantees for every concrete function, and they are not a
claim that no structural cryptanalysis can do better. This chapter makes no
quantum work-factor claim.

\section{Diffusion Is Not a Security Proof}

The draft called the avalanche effect a security property. It is better treated
here as an empirical diffusion observation: a small input change often changes
many digest bits.

The companion hashes two 32-byte messages that differ in one input bit. Their
BLAKE2b-256 digests differ in 130 of 256 positions. This is visually useful,
but it proves neither preimage resistance, second-preimage resistance, nor
collision resistance. A deliberately weak construction can be made to look
well diffused on selected examples.

\section{What an Unkeyed Digest Can Establish}

An unkeyed digest is useful as a reproducible fingerprint of exact bytes. A
mismatch against an authentic reference detects that the bytes differ. Digest
equality, however, is conditional evidence: collisions exist. Against an
adversary, the needed assumption depends on the experiment---second-preimage
resistance when one fixed original is already bound, or collision resistance
when an attacker can prepare both alternatives. An attacker who can replace
both a file and its reference digest can always make the replacement internally
consistent.

\begin{center}
\small
\begin{tabular}{|>{\raggedright\arraybackslash}p{0.25\textwidth}|>{\centering\arraybackslash}p{0.14\textwidth}|>{\raggedright\arraybackslash}p{0.45\textwidth}|}
\hline
Goal & Bare digest? & What is additionally required \\
\hline
Reproducible fingerprint & Yes & Exact bytes, algorithm, parameters, and output length \\
\hline
Detect change & Conditional & Authentic reference, exact encoding, and the relevant resistance assumption \\
\hline
Authenticate a sender & No & A specified MAC or signature construction and keys \\
\hline
Hide a guessable value & No & A defined hiding construction and its assumptions \\
\hline
Store a password & No & A salted, costed password-hashing scheme \\
\hline
Derive key material & Not generally & A specified KDF or extractor and entropy assumptions \\
\hline
Provide anonymity & No & An anonymity system with a leakage and adversary model \\
\hline
\end{tabular}
\end{center}

\subsection{Low-entropy inputs are guessable}

Publishing \(H(\mathtt{paper})\) does not hide an RPS move. An observer hashes
the three public candidates \(\mathtt{rock}\), \(\mathtt{paper}\), and
\(\mathtt{scissors}\), then compares the results. The C demonstration performs
this complete dictionary attack and recovers \(\mathtt{paper}\).

The commitment in the book's Rock--Paper--Scissors chapter therefore hashes a
canonical context, the move, and a fresh secret high-entropy nonce. Even that
larger construction retains explicitly conditional hiding and binding
arguments; the bare hash alone supplies neither property.

\section{Adjacent Constructions Are Separate Protocols}

\paragraph{Message authentication.}
A MAC includes a secret key and a specified construction. HMAC, for example,
is not the operation \(H(m)\); it is a keyed construction defined separately
\cite{rfc2104hmac}. BLAKE2 also has a keyed mode, but this appendix's wrapper
deliberately exposes only unkeyed hashing.

\paragraph{Digital signatures.}
A signature scheme defines key generation, signing, and verification. Many
schemes hash an encoded message internally, but a digest by itself identifies
no signer and provides no signature.

\paragraph{Commitments.}
A commitment needs both a hiding objective and a binding objective, plus a
defined opening procedure and context. A bare digest of a predictable value is
not hiding. The project's educational commitment is catalogued separately in
the primitive registry.

\paragraph{Password hashing.}
Fast general-purpose hashing makes offline guessing cheap. Password storage
needs a dedicated salted scheme with tunable resource costs. The official
libsodium documentation explicitly says that its BLAKE2b generic-hash API is
not suitable for password hashing \cite{libsodium-generichash}.

\paragraph{Key derivation.}
Hash functions can be components of KDFs and extractors, but a defined KDF does
more than shorten bytes. Hashing weak or predictable material does not create
entropy.

\paragraph{Confidentiality and anonymity.}
An unkeyed digest is deterministic and public. It is not encryption, and it
can create stable linkable identifiers. Confidentiality and anonymity require
their own constructions, threat models, and leakage analyses.

\section{Structured Inputs and Domain Separation}

Raw concatenation is ambiguous:
\[
    \mathtt{ab} \mathbin\| \mathtt{c}
    =
    \mathtt{a} \mathbin\| \mathtt{bc}
    =
    \mathtt{abc}.
\]
This equality is not a hash collision; both sides are literally the same byte
string. A protocol must encode boundaries before hashing, for example with
fixed-width fields or unambiguous lengths. Standard derived functions such as
TupleHash likewise encode tuple elements so different tuples remain different
inputs \cite{nist800185}.

A public domain label and version can also prevent the same bytes from being
silently interpreted as different protocol objects. Domain separators are
context, not secrets. They do not replace a MAC key, encryption key, nonce, or
password salt.

\section{The Project's BLAKE2b-256 Profile}

The companion fixes one profile:

\begin{itemize}
    \item algorithm: unkeyed BLAKE2b;
    \item output length: 32 bytes;
    \item implementation backend: libsodium; and
    \item input: the exact caller-supplied byte string and length.
\end{itemize}

RFC~7693 specifies BLAKE2b and lists the 32-byte parameter set with a
\(2^{128}\) collision-security target label \cite{saarinen2015blake2}. The same
RFC explicitly makes no independent assertion of BLAKE2's security. The table
label is therefore recorded as a design target, not as a proof.

The output length is part of BLAKE2 initialization. BLAKE2b-256 is not defined
as the first 32 bytes of a BLAKE2b-512 digest. Algorithm name, parameters, and
exact input bytes all belong beside a published digest.

\section{MD5: A Historical Collision Failure}

The old draft intended to demonstrate an MD5 collision. That instinct is worth
preserving, but the recorded blocks contained corrupted nibbles and did not
collide. The active chapter replaces them with a verified historical vector in
\path{test-vectors/hash-foundations-v1.md}; the untouched backup preserves the
original draft.

The short examples are:

\begin{verbatim}
MD5("Hello, World!") = 65a8e27d8879283831b664bd8b7f0ad4
MD5("ello, World!")  = 2f695452b98fa1028dc11e1e60f17614
MD5(empty)           = d41d8cd98f00b204e9800998ecf8427e
\end{verbatim}

The former chapter incorrectly printed the empty-input digest for
\texttt{ello, World!}. Correcting a test vector does not rehabilitate MD5.
RFC~6151 concludes that MD5 is no longer acceptable where collision resistance
is required, including digital signatures \cite{rfc6151md5}. The verified
128-byte pair in the vector file consists of distinct messages with the same
MD5 digest, illustrating the collision task studied by Wang and Yu
\cite{wang2005md5}.

MD5 is not compiled or linked into this project's C companion. It remains only
an attack-history exercise. The example also teaches a useful distinction: a
collision break does not mean that every preimage problem has become equally
easy.

\section{C Implementation Companion}

The public interface is \path{include/ac/hash.h}; the implementation is
\path{src/primitives/hash.c}; and the teaching driver is
\path{src/hash_functions/hash_demo.c}.

The public operation is intentionally narrow:

\begin{verbatim}
ac_status ac_hash_blake2b_256(
    ac_hash_blake2b_256_digest *digest,
    const uint8_t *input,
    size_t input_len);
\end{verbatim}

It validates the null/length policy, pins the algorithm and output length,
checks library initialization and return values, and clears a non-null output
on detected failure. It does not expose a generic framework, keyed mode,
password API, or libsodium state type. The input and output buffers must not
overlap.

Run the fixed public demonstration with:

\begin{verbatim}
cd /src
make -C demostrations demo_hash
\end{verbatim}

The demonstration prints an explicit warning, the empty and \texttt{abc}
vectors, the one-bit diffusion observation, the complete three-choice guessing
attack, and the ambiguous-concatenation lesson.

\section{Tests and Evidence Boundary}

The deterministic vectors live in
\path{test-vectors/hash-foundations-v1.md}. The C test in
\path{tests/test_hash.c} checks:

\begin{itemize}
    \item BLAKE2b-256 of the empty input and \texttt{abc};
    \item all byte values from \texttt{00} through \texttt{ff};
    \item binary lengths, including \texttt{abc} versus \texttt{abc\textbackslash0};
    \item determinism and the exact 130-bit diffusion observation;
    \item recovery of a bare-hashed three-choice input; and
    \item invalid-argument handling and output clearing.
\end{itemize}

The MD5 collision was independently reproduced with a standard-library tool
that is not a project runtime dependency. The two messages also have distinct
BLAKE2b-256 digests, as recorded in the vector file.

Passing these tests establishes exact behavior for the recorded implementation
and inputs. It does not establish computational hardness, resistance to an
unknown attack, production suitability, or the security of a larger protocol.

\section{What Cryptography Does Not Solve}

A digest records a deterministic relationship between bytes. It does not
determine whether those bytes are truthful, whether their author is
trustworthy, or whether their interpretation is just. An unkeyed digest does
not authenticate its own reference value, conceal a guessable input, repair
ambiguous serialization, protect a compromised endpoint, force a party to
reveal a commitment, or guarantee the security of a larger protocol. Those
properties require additional constructions and assumptions.

\section{Limitations and Open Questions}

This appendix does not formalize all seven notions catalogued by Rogaway and
Shrimpton, survey BLAKE2 cryptanalysis, analyze quantum queries, benchmark
implementations, or approve an algorithm for production deployment. It gives a
fixed educational profile and the minimum distinctions required by later
chapters. Each later construction must state which property it needs and why
its complete encoding and threat model justify that use.

\section{Exercises}

\begin{enumerate}
    \setlength{\itemsep}{0.25em}
    \item Hash \texttt{abc}, \texttt{abc\textbackslash n}, and
    \texttt{abc\textbackslash0}; record the exact input bytes.
    \item Explain why publishing \(H(\mathtt{rock})\) does not hide an RPS move.
    \item Construct an active replacement attack when a file and its digest
    share an unauthenticated channel.
    \item Derive the approximate collision probability after \(q\) ideal
    \(n\)-bit queries.
    \item Explain why a collision pair need not solve a chosen preimage.
    \item Repair \(\mathtt{a}\|\mathtt{bc}\) versus
    \(\mathtt{ab}\|\mathtt{c}\) with an explicit canonical encoding.
    \item Reproduce the historical MD5 pair; state which property it falsifies.
    \item Reproduce the commitment vector; distinguish byte correctness from
    conditional hiding and binding.
\end{enumerate}
\flushbottom
